\begin{exercice*}[Extrait du DNB]
    Sam préfère les bonbons bleus. Dans son paquet de 500 bonbons, 150 sont bleus, les autres sont rouges, jaunes ou verts.

    \begin{enumerate}
        \item Quelle est la probabilité qu'il pioche au hasard un bonbon bleu dans son paquet ?
        \item \SI{20}{\percent} des bonbons de ce paquet sont rouges. Combien de bonbons rouges y a-t-il ?
        \item Sachant qu'il y a 130 bonbons verts dans ce paquet, Sam a-t-il plus de chance de piocher au hasard un bonbon vert ou un bonbon jaune ?
    \end{enumerate}
\end{exercice*}
\begin{corrige}
    %\setcounter{partie}{0} % Pour s'assurer que le compteur de \partie est à zéro dans les corrigés
    %\phantom{rrr}    
    \begin{enumerate}
        \item La probabilité qu'il pioche au hasard un bonbon bleu dans son paquet est de $\dfrac{150}{500} = \num{0,3}$.
        \item Il y a \SI{20}{\percent} des bonbons de ce paquet qui sont rouges. Donc il y a $500 \times \num{0,2} = 100$ bonbons rouges.
        \item Il y a 130 bonbons verts dans ce paquet. Il reste donc $500 - 150 - 100 - 130 = 120$ bonbons jaunes. Sam a donc plus de chance de piocher un bonbon jaune que un bonbon vert.
    \end{enumerate}
\end{corrige}

